% \documentclass[polish, 8pt]{article}
\documentclass[polish, 10pt]{extarticle}
% \usepackage[top=2cm, bottom=2cm, left=3cm, right=3cm]{geometry}
\usepackage[a4paper,top=1.2cm,bottom=1.2cm,left=1.2cm,right=1.2cm]{geometry}
\usepackage{graphicx} % Required for inserting images

\usepackage[polish]{babel}
\usepackage{polski}
\usepackage{placeins}
\usepackage{dirtytalk}
% Set page size and margins
% Replace `letterpaper' with `a4paper' for UK/EU standard size
\usepackage{csquotes}
% Useful packages
\usepackage{float}
\usepackage{amsmath}
\usepackage{graphicx}
\usepackage{hyperref}
\usepackage{graphicx}
\usepackage{listings}
\usepackage{xcolor}
\usepackage{amssymb}
\usepackage{enumitem}
\usepackage{float}
\definecolor{codegreen}{rgb}{0,0.6,0}
\definecolor{codegray}{rgb}{0.5,0.5,0.5}
\definecolor{codepurple}{rgb}{0.58,0,0.82}
\definecolor{backcolour}{rgb}{0.95,0.95,0.92}

\lstdefinestyle{mystyle}{
    backgroundcolor=\color{backcolour},   
    commentstyle=\color{codegreen},
    keywordstyle=\color{magenta},
    numberstyle=\tiny\color{codegray},
    stringstyle=\color{codepurple},
    basicstyle=\ttfamily\footnotesize,
    breakatwhitespace=false,         
    breaklines=true,                 
    captionpos=b,                    
    keepspaces=true,                 
    numbers=left,                    
    numbersep=5pt,                  
    showspaces=false,                
    showstringspaces=false,
    showtabs=false,                  
    tabsize=2
}
\lstset{style=mystyle}
\hypersetup{
    colorlinks=true,
    linktoc=all     %set to all if you want both sections and subsections linked
}
\usepackage{mdframed}
\usepackage{titling}
\setlength{\droptitle}{-2cm} 

\begin{document}

\title{Kolokwium II -- Zestaw B}
\author{}
\date{}

\maketitle
\vspace{-2cm}

\section*{Zadanie 1 (30 pkt)}
Niech dwuwymiarowa zmienna losowa $(X, Y)$ ma rozkład ciągły o gęstości:
\[
f(x,y) = 
\begin{cases} 
cx & \text{dla } 0 \le x \le 1, \,\, x^2 \le y \le 1, \\
0 & \text{w przeciwnym przypadku,}
\end{cases}
\]
gdzie $c \in \mathbb{R_+}$ jest pewną stałą.

\begin{enumerate}
    \item[a)] \textbf{Wyznacz wartość stałej $c$} \hfill (5 pkt).
    \item[b)] \textbf{Wyznacz gęstości brzegowe} \hfill (5 pkt).
    \item[c)] \textbf{Wyznacz gęstości warunkowe} \hfill (5 pkt).
    \item[d)] \textbf{Wyznacz $\mathbb P\!\left(Y>\tfrac12\mid X=\tfrac12\right)$} \hfill (5 pkt).
    \item[e)] \textbf{Oblicz kowariancję $\text{Cov}(X, Y)$} \hfill (5 pkt).
    \item[f)] \textbf{Opisz sposób wyznaczenia współczynnika korelacji $\rho_{X,Y}$},\\rozpisując wzory i całki niezbędne do jego obliczenia \hfill (5 pkt).
\end{enumerate}

\section*{Rozwiązanie -- Zadanie 1}

Dany jest rozkład ciągły $(X,Y)$ o gęstości
\[
f(x,y)=
\begin{cases}
cx & \text{dla } 0\le x\le 1,\;\; x^2\le y\le 1,\\
0 & \text{w p.p.}
\end{cases}
\]
Obszar, gdzie $f>0$ ma postać
\[
D=\{(x,y): 0\le x\le 1,\; x^2\le y\le 1\}.
\]

\begin{enumerate}
\item[a)] \textbf{Wyznaczenie stałej $c$.}

Warunek normalizacji:
\[
\iint_{\mathbb{R}^2} f(x,y)\,dx\,dy=1
\quad\Longleftrightarrow\quad
\int_{0}^{1}\int_{x^2}^{1} cx\,dy\,dx=1.
\]
Liczymy całki:
\[
\int_{0}^{1}\int_{x^2}^{1} cx\,dy\,dx
=\int_{0}^{1} cx\,(1-x^2)\,dx
=c\int_{0}^{1}(x-x^3)\,dx
=c\left(\frac12-\frac14\right)=\frac{c}{4}.
\]
Zatem $\frac{c}{4}=1$, czyli
\[
\boxed{c=4.}
\]

\item[b)] \textbf{Gęstości brzegowe.}

\underline{Uwaga (częsty błąd!)} Przy wyznaczaniu gęstości brzegowych trzeba \textbf{koniecznie podać poprawne przedziały}, tzn. dla jakich $x$ (lub $y$) dana gęstość jest dodatnia. Za brak przedziałów odejmowałem 1 punkt.

\smallskip
\textbf{Gęstość brzegowa $X$.}
Dla ustalonego $x$ gęstość jest niezerowa wtedy i tylko wtedy, gdy $0\le x\le 1$ oraz $y\in[x^2,1]$. Zatem
\[
f_X(x)=\int_{-\infty}^{\infty} f(x,y)\,dy=\int_{x^2}^{1} 4x\,dy
=4x(1-x^2),
\]
czyli
\[
\boxed{
f_X(x)=
\begin{cases}
4x(1-x^2), & 0\le x\le 1,\\
0, & \text{w p.p.}
\end{cases}}
\]

\smallskip
\textbf{Gęstość brzegowa $Y$.}
Dla ustalonego $y$ warunki $0\le x\le 1$ i $x^2\le y\le 1$ oznaczają:
\[
0\le y\le 1
\quad\text{oraz}\quad
0\le x\le \sqrt{y}.
\]
Stąd
\[
f_Y(y)=\int_{-\infty}^{\infty} f(x,y)\,dx=\int_{0}^{\sqrt{y}} 4x\,dx
=4\cdot \frac{x^2}{2}\Big|_{0}^{\sqrt{y}}=2y,
\]
czyli
\[
\boxed{
f_Y(y)=
\begin{cases}
2y, & 0\le y\le 1,\\
0, & \text{w p.p.}
\end{cases}}
\]

\item[c)] \textbf{Gęstości warunkowe.}

\underline{Uwaga (częsty błąd!)} Tak jak w punkcie (b), \textbf{trzeba podać warunki/przedziały dla obu zmiennych}.

\smallskip
\textbf{Gęstość $Y|X=x$.}
Dla $0<x<1$ mamy $f_X(x)=4x(1-x^2)>0$ oraz
\[
f_{Y|X=x}(y)=\frac{f(x,y)}{f_X(x)}
=\frac{4x}{4x(1-x^2)}=\frac{1}{1-x^2} \quad\text{dla } y\in[x^2,1].
\]

\smallskip
\textbf{Gęstość $X|Y=y$.}
Dla $0<y\le 1$ mamy $f_Y(y)=2y>0$. Ponadto nośnik w $x$ to $0\le x\le \sqrt{y}$. Zatem
\[
f_{X|Y=y}(x)=\frac{f(x,y)}{f_Y(y)}=\frac{4x}{2y}=\frac{2x}{y}
\quad\text{dla } 0\le x\le \sqrt{y}.
\]

\item[d)] \textbf{Obliczenie $\mathbb{P}\!\left(Y>\frac12\mid X=\frac12\right)$.}

Najpierw ustalamy \textbf{nośnik warunkowy} $Y|X=\tfrac12$.
Jeśli $X=\tfrac12$, to z warunku nośnika $x^2\le y\le 1$ dostajemy
\[
\left(\frac12\right)^2\le y\le 1 \quad\Longleftrightarrow\quad \frac14\le y\le 1.
\]
\underline{Uwaga (częsty błąd!)} Trzeba uwzględnić dla jakiego przedziału gęstość warunkowa jest niezerowa i w ten sposób uzasadnić dobrane granie w całce. Za brak uzasadnienia odejmowałem 1 punkt.

Korzystamy z gęstości warunkowej z (c):
\[
f_{Y|X=1/2}(y)=\frac{1}{1-(1/2)^2}=\frac{1}{1-\frac14}=\frac{4}{3},
\quad y\in\left[\frac14,1\right].
\]
Zatem
\[
\mathbb{P}\!\left(Y>\frac12\mid X=\frac12\right)
=\int_{1/2}^{1} \frac{4}{3}\,dy
=\frac{4}{3}\left(1-\frac12\right)=\frac{2}{3}.
\]
Ostatecznie
\[
\boxed{\mathbb{P}\!\left(Y>\frac12\mid X=\frac12\right)=\frac{2}{3}.}
\]

\item[e)] \textbf{Kowariancja $\mathrm{Cov}(X,Y)$.}

\[
\mathrm{Cov}(X,Y)=\mathbb{E}[XY]-\mathbb{E}[X]\mathbb{E}[Y].
\]

\smallskip
\textbf{1) Liczymy $\mathbb{E}[X]$.}
\[
\mathbb{E}[X]=\int_0^1\int_{x^2}^{1} x\cdot 4x\,dy\,dx
=\int_0^1 4x^2(1-x^2)\,dx
=4\int_0^1 (x^2-x^4)\,dx
=4\left(\frac13-\frac15\right)=\frac{8}{15}.
\]

\smallskip
\textbf{2) Liczymy $\mathbb{E}[Y]$.} Wygodnie z gęstości brzegowej:
\[
\mathbb{E}[Y]=\int_0^1 y\,f_Y(y)\,dy=\int_0^1 y\cdot 2y\,dy
=2\int_0^1 y^2\,dy=\frac{2}{3}.
\]

\smallskip
\textbf{3) Liczymy $\mathbb{E}[XY]$.}
\[
\mathbb{E}[XY]=\int_0^1\int_{x^2}^{1} xy\cdot 4x\,dy\,dx
=\int_0^1 4x^2\left(\int_{x^2}^{1} y\,dy\right)\,dx.
\]
\[
\int_{x^2}^{1} y\,dy=\frac{y^2}{2}\Big|_{x^2}^{1}=\frac{1-x^4}{2},
\]
więc
\[
\mathbb{E}[XY]=\int_0^1 4x^2\cdot \frac{1-x^4}{2}\,dx
=2\int_0^1 (x^2-x^6)\,dx
=2\left(\frac13-\frac17\right)=\frac{8}{21}.
\]

\smallskip
Zatem
\[
\mathrm{Cov}(X,Y)=\frac{8}{21}-\frac{8}{15}\cdot\frac{2}{3}
=\frac{8}{21}-\frac{16}{45}
=\frac{120}{315}-\frac{112}{315}
=\frac{8}{315}.
\]
\[
\boxed{\mathrm{Cov}(X,Y)=\frac{8}{315}.}
\]

\item[f)] \textbf{Opis sposobu wyznaczenia współczynnika korelacji $\rho_{X,Y}$.}

Współczynnik korelacji:
\[
\rho_{X,Y}=\frac{\mathrm{Cov}(X,Y)}{\sqrt{\mathrm{Var}(X)\,\mathrm{Var}(Y)}}.
\]
A więc potrzebujemy:
\[
\mathrm{Cov}(X,Y)=\mathbb{E}[XY]-\mathbb{E}[X]\mathbb{E}[Y],
\quad
\mathrm{Var}(X)=\mathbb{E}[X^2]-(\mathbb{E}[X])^2,
\quad
\mathrm{Var}(Y)=\mathbb{E}[Y^2]-(\mathbb{E}[Y])^2.
\]

\smallskip
\textbf{Krok 1: policzyć wartości oczekiwane} (podpunkt e):
\[
\mathbb{E}[X]=\int_{0}^{1}\int_{x^2}^{1} x\,f(x,y)\,dy\,dx,
\qquad
\mathbb{E}[Y]=\int_{0}^{1}\int_{x^2}^{1} y\,f(x,y)\,dy\,dx,
\]
\[
\mathbb{E}[XY]=\int_{0}^{1}\int_{x^2}^{1} xy\,f(x,y)\,dy\,dx.
\]

\smallskip
\textbf{Krok 2: policzyć drugie momenty}:
\[
\mathbb{E}[X^2]=\int_{0}^{1}\int_{x^2}^{1} x^2\,f(x,y)\,dy\,dx
\quad\text{(również można: } \mathbb{E}[X^2]=\int_0^1 x^2 f_X(x)\,dx\text{)},
\]
\[
\mathbb{E}[Y^2]=\int_{0}^{1}\int_{x^2}^{1} y^2\,f(x,y)\,dy\,dx
\quad\text{(również można: } \mathbb{E}[Y^2]=\int_0^1 y^2 f_Y(y)\,dy\text{)}.
\]

\smallskip
\textbf{Krok 3: złożyć wszystko we wzór na korelację}:
\[
\rho_{X,Y}
=\frac{\mathbb{E}[XY]-\mathbb{E}[X]\mathbb{E}[Y]}
{\sqrt{\big(\mathbb{E}[X^2]-(\mathbb{E}[X])^2\big)\big(\mathbb{E}[Y^2]-(\mathbb{E}[Y])^2\big)}}.
\]

\end{enumerate}

\section*{Zadanie 2 (15 pkt)}
Dla zmiennej losowej $X\sim \mathrm{Pois}(\lambda)$:
\begin{enumerate}
  \item[a)] Wyznacz funkcję tworzącą momenty (MGF) \hfill (10 pkt).
  \item[b)] Korzystając z MGF oblicz wariancję $\mathrm{Var}(X)$ \hfill (5 pkt).
\end{enumerate}

\section*{Rozwiązanie -- Zadanie 2}

Dla zmiennej losowej $X\sim \mathrm{Pois}(\lambda)$ mamy
\[
\mathbb{P}(X=k)=e^{-\lambda}\frac{\lambda^k}{k!},\quad k=0,1,2,\dots
\]

\begin{enumerate}
\item[a)] \textbf{Wyznaczenie funkcji tworzącej momenty (MGF).}

Z definicji:
\[
M_X(t)=\mathbb{E}\!\left(e^{tX}\right)=\sum_{k=0}^{\infty} e^{tk}\,\mathbb{P}(X=k).
\]
Podstawiamy postać $\mathbb{P}(X=k)$:
\[
\begin{aligned}
M_X(t)
&=\sum_{k=0}^{\infty} e^{tk}\,e^{-\lambda}\frac{\lambda^k}{k!}
= e^{-\lambda}\sum_{k=0}^{\infty}\frac{(\lambda e^t)^k}{k!}
= e^{-\lambda}\exp(\lambda e^t)\\
&=\exp\!\big(\lambda(e^t-1)\big).
\end{aligned}
\]
Wyżej korzystamy z rozwinięcia wykładniczej w szereg potęgowy:
\[
e^u=\sum_{k=0}^{\infty}\frac{u^k}{k!}\quad \text{dla każdego } u\in\mathbb{R}.
\]
Zatem
\[
\boxed{M_X(t)=\exp\!\big(\lambda(e^t-1)\big).}
\]

\item[b)] \textbf{Obliczenie wariancji $\mathrm{Var}(X)$ z MGF.}

\smallskip

Najpierw liczymy wartość oczekiwaną:
\[
\mathbb{E}(X)=M_X'(0).
\]
Liczymy pochodną:
\[
M_X(t)=\exp\!\big(\lambda(e^t-1)\big)
\quad\Longrightarrow\quad
M_X'(t)=\exp\!\big(\lambda(e^t-1)\big)\cdot \lambda e^t
=\lambda e^t\,M_X(t).
\]
Podstawiamy $t=0$:
\[
M_X'(0)=\lambda e^0\,M_X(0)=\lambda\cdot 1\cdot \exp(\lambda(1-1))=\lambda.
\]
Stąd:
\[
\mathbb{E}(X)=\lambda.
\]

Do wyznaczenia wariancji korzystamy ze wzoru:

\[
\boxed{\mathrm{Var}(X)=\mathbb{E}(X^2)-(\mathbb{E}(X))^2}
\]

Z własności MGF:
\[
\mathbb{E}(X^2)=M_X''(0).
\]
Mamy już
\[
M_X'(t)=\lambda e^t\,M_X(t).
\]
Różniczkujemy drugi raz:
\[
\begin{aligned}
M_X''(t)
&=\frac{d}{dt}\big(\lambda e^t M_X(t)\big)
=\lambda e^t M_X(t)+\lambda e^t M_X'(t).
\end{aligned}
\]
Podstawiamy ponownie $M_X'(t)=\lambda e^t M_X(t)$:
\[
M_X''(t)=\lambda e^t M_X(t)+\lambda e^t\big(\lambda e^t M_X(t)\big)
=\big(\lambda e^t+\lambda^2 e^{2t}\big)M_X(t).
\]
Podstawiamy $t=0$:
\[
M_X''(0)=\big(\lambda+\lambda^2\big)M_X(0)=\lambda+\lambda^2.
\]
Zatem
\[
\mathbb{E}(X^2)=\lambda+\lambda^2,\qquad \mathbb{E}(X)=\lambda,
\]
a więc
\[
\mathrm{Var}(X)=\mathbb{E}(X^2)-(\mathbb{E}(X))^2
=(\lambda+\lambda^2)-\lambda^2=\lambda.
\]
Ostatecznie
\[
\boxed{\mathrm{Var}(X)=\lambda.}
\]

\end{enumerate}

\section*{Zadanie 3 (30 pkt)}
System zasilania awaryjnego w serwerowni składa się z $n$ niezależnych modułów bateryjnych. Moduły te włączają się sekwencyjnie – gdy pierwszy się wyczerpie, natychmiast startuje drugi, itd. Czas pracy każdego $i$-tego modułu (w godzinach) jest zmienną losową $X_i$ o rozkładzie wykładniczym z parametrem $\lambda = 2$: $X_i \sim \text{Exp}(2)$. Niech $T$ oznacza całkowity czas pracy systemu do momentu wyczerpania wszystkich baterii:
\[
T = X_1 + X_2 + \dots + X_n
\]

\begin{enumerate}
    \item[a)] Oblicz wartość oczekiwaną $\mathbb{E}[T]$ i wariancję $\text{Var}(T)$ \hfill (5 pkt).
    \item[b)] Korzystając z własności sumy niezależnych zmiennych losowych, podaj dokładny rozkład zmiennej $T$ \hfill (5 pkt).
    \item[c)] Załóżmy, że system składa się z $n=100$ modułów. Oszacuj prawdopodobieństwo, że system będzie działał dłużej niż 55 godzin. Powołaj się na odpowiednie twierdzenie \hfill (10 pkt).
    \item[d)] Traktując parametr $\lambda$ jako nieznany, a czasy pracy modułów $x_1,\dots,x_n$ jako zaobserwowane realizacje niezależnych zmiennych $X_1,\dots,X_n \sim \mathrm{Exp}(\lambda)$, wyznacz estymator największej wiarygodności $\hat\lambda_{\text{MLE}}$. \hfill (10 pkt)
\end{enumerate}

\section*{Rozwiązanie -- Zadanie 3}

Mamy $n$ niezależnych modułów, a czas pracy $i$-tego modułu:
\[
X_i\sim \mathrm{Exp}(2),\qquad f_{X}(x)=2e^{-2x}\ (x\ge 0).
\]
Całkowity czas pracy systemu:
\[
T=X_1+\cdots+X_n.
\]

\begin{enumerate}
\item[a)] \textbf{Wartość oczekiwana i wariancja $T$.}

Dla rozkładu wykładniczego $X\sim \mathrm{Exp}(\lambda)$ zachodzi:
\[
\mathbb{E}[X]=\frac{1}{\lambda},\qquad \mathrm{Var}(X)=\frac{1}{\lambda^2}.
\]
Stąd dla $\lambda=2$:
\[
\mathbb{E}[X_i]=\frac{1}{2},\qquad \mathrm{Var}(X_i)=\frac{1}{4}.
\]
Ponieważ $X_1,\dots,X_n$ są \textbf{niezależne}:
\[
\mathbb{E}[T]=\sum_{i=1}^n \mathbb{E}[X_i]=n\cdot \frac{1}{2}=\frac{n}{2},
\]
\[
\mathrm{Var}(T)=\sum_{i=1}^n \mathrm{Var}(X_i)=n\cdot \frac{1}{4}=\frac{n}{4}.
\]
\[
\boxed{\mathbb{E}[T]=\frac{n}{2},\qquad \mathrm{Var}(T)=\frac{n}{4}.}
\]

\item[b)] \textbf{Dokładny rozkład zmiennej $T$.}

Najpierw zauważmy, że rozkład wykładniczy jest szczególnym przypadkiem rozkładu gamma:
\[
X\sim \mathrm{Exp}(\lambda)\quad \Longleftrightarrow \quad X\sim Gamma(1,\ \lambda),
\]
Zatem dla każdego $i$:
\[
X_i\sim Gamma(1,2).
\]

\textbf{Własność (suma niezależnych zmiennych gamma o tym samym parametrze rate).}
Jeżeli $Y_1,\dots,Y_n$ są niezależne oraz
\[
Y_i\sim Gamma(s_i,\lambda)\quad \text{dla tego samego }\lambda,
\]
to
\[
\sum_{i=1}^n Y_i \sim Gamma\!\left(\sum_{i=1}^n s_i,\ \lambda\right).
\]

Stosując to do $T=\sum_{i=1}^n X_i$ z $s_i=1$ i $\lambda=2$:
\[
T\sim \Gamma\!\left(\sum_{i=1}^n 1,\ 2\right)=\Gamma(n,2).
\]

Gęstość $T$:
\[
f_T(t)=\frac{2^n}{\Gamma(n)}t^{n-1}e^{-2t},\qquad t> 0.
\]
\[
\boxed{T\sim Gamma(n,2)}
\]

\item[c)] \textbf{Dla $n=100$: oszacowanie $\mathbb{P}(T>55)$.}

Z punktu (a):
\[
\mu=\mathbb{E}[T]=\frac{100}{2}=50,\qquad \sigma^2=\mathrm{Var}(T)=\frac{100}{4}=25,\qquad \sigma=5.
\]

\textbf{Twierdzenie (Centralne Twierdzenie Graniczne).}
Dla dużego $n$ zmienna
\[
\frac{T-\mu}{\sigma}\approx \mathcal{N}(0,1).
\]
Zatem:
\[
\mathbb{P}(T>55)=\mathbb{P}\!\left(\frac{T-50}{5}>\frac{55-50}{5}\right)
\approx \mathbb{P}(Z>1),
\]
gdzie $Z\sim \mathcal{N}(0,1)$.
\[
\boxed{\mathbb{P}(Z>1)=1-\Phi(1)}
\]

\item[d)] \textbf{Estymator największej wiarygodności $\hat\lambda_{\mathrm{MLE}}$ dla $\lambda$ przy $X_i\sim \mathrm{Exp}(\lambda)$.}

Załóżmy obserwacje $x_1,\dots,x_n$ niezależne oraz
\[
f(x\mid \lambda)=\lambda e^{-\lambda x}\ (x\ge 0),\qquad \lambda>0.
\]

\textbf{Krok 1: funkcja wiarygodności.}
\[
L(\lambda)=\prod_{i=1}^n \lambda e^{-\lambda x_i}
=\lambda^n \exp\!\left(-\lambda\sum_{i=1}^n x_i\right).
\]

\textbf{Krok 2: log-wiarygodność.}
\[
\ell(\lambda)=\ln L(\lambda)=n\ln\lambda-\lambda\sum_{i=1}^n x_i.
\]

\textbf{Krok 3: pochodna i warunek konieczny.}
\[
\ell'(\lambda)=\frac{n}{\lambda}-\sum_{i=1}^n x_i.
\]
Równanie $\ell'(\lambda)=0$ daje:
\[
\frac{n}{\lambda}=\sum_{i=1}^n x_i
\quad\Longrightarrow\quad
\hat\lambda_{\mathrm{MLE}}=\frac{n}{\sum_{i=1}^n x_i}.
\]

\textbf{Krok 4: sprawdzenie maksimum.}
\[
\ell''(\lambda)=-\frac{n}{\lambda^2}<0\quad (\lambda>0),
\]
więc otrzymany punkt jest maksimum.

\[
\boxed{\hat\lambda_{\mathrm{MLE}}=\dfrac{n}{\sum_{i=1}^n x_i}=\dfrac{1}{\bar x}.}
\]

\end{enumerate}

\section*{Zadanie 4 (25 pkt)}
W pewnym zakładzie produkcyjnym kontroluje się czas (w minutach) potrzebny na wykonanie standardowej operacji. Przyjmuje się, że czasy wykonywania tej operacji w populacji mają rozkład normalny. Wylosowano próbę \textbf{$n = 14$} pomiarów, dla których uzyskano $\bar{x} = 12{,}8$ min przy $s = 3{,}2$ min.

\begin{enumerate}
    \item \textbf{Opisz sposób konstrukcji $95\%$ przedziału ufności dla $\mu$}, tj.:
    \begin{itemize}
        \item jaka statystyka jest wykorzystywana do budowy przedziału i jaki jest jej rozkład, \hfill (5 pkt),
        \item dlaczego wybieramy akurat tę statystykę w tej sytuacji, \hfill (5 pkt),
        \item opisz sposób konstrukcji (jeśli wprowadzasz zmienne, wyjaśnij  co oznaczają) \hfill (10 pkt).
    \end{itemize}

    \item Ktoś interpretuje wynik następująco:
    \begin{quote}
    ``Z prawdopodobieństwem $95\%$ prawdziwa wartość średniej $\mu$ leży w wyznaczonym przedziale.''
    \end{quote}
    Czy ta interpretacja jest \textbf{prawdziwa}? Uzasadnij. \hfill (5 pkt).
\end{enumerate}

\section*{Rozwiązanie -- Zadanie 4}

\begin{enumerate}
\item \textbf{Opis sposobu konstrukcji $95\%$ przedziału ufności dla $\mu$.}

\begin{itemize}
\item[a)] \textbf{Jaka statystyka jest wykorzystywana do budowy przedziału, jaki jest jej rozkład.}

Ponieważ $\sigma$ jest nieznane, wykorzystujemy statystykę t-Studenta:
\[
T=\frac{\bar{X}-\mu}{S/\sqrt{n}},
\]
gdzie $\bar{X}$ jest średnią z próby, a $S$ jest odchyleniem standardowym z próby.
Przy założeniu normalności populacji mamy
\[
T\sim t_{n-1}.
\]

\item[b)] \textbf{Dlaczego wybieramy akurat tę statystykę w tej sytuacji?}

Ponieważ zakładamy, że populacja ma rozkład normalny, to średnia z próby $\bar X$ jest naturalnym estymatorem $\mu$.
Kluczowe jest jednak to, że nie znamy odchylenia standardowego populacji $\sigma$, więc w standaryzacji musimy je zastąpić
odchyleniem z próby $S$. Wtedy statystyka
\[
T=\frac{\bar{X}-\mu}{S/\sqrt{n}}
\]
ma \emph{dokładnie} rozkład t-Studenta z $n-1$ stopniami swobody (przy założeniu normalności populacji). Dlatego
przedział ufności dla $\mu$ konstruujemy w oparciu o kwantyle rozkładu $t_{n-1}$.

Gdyby $\sigma$ było znane, to moglibyśmy użyć statystyki
\[
Z=\frac{\bar{X}-\mu}{\sigma/\sqrt{n}}\sim\mathcal{N}(0,1),
\]
ale w tym zadaniu $\sigma$ jest nieznane, więc poprawnym wyborem jest $t$-statystyka.

\item[c)] \textbf{Krok 1: ``wycięcie'' środkowych $95\%$ rozkładu statystyki.}

Ponieważ
\[
T\sim t_{n-1},
\]
to dla $\alpha=0{,}05$ zachodzi
\[
\mathbb{P}\!\left(t_{\alpha/2,\;n-1}\le T\le t_{1-\alpha/2,\;n-1}\right)=1-\alpha=0{,}95,
\]
czyli wprost
\[
\mathbb{P}\!\left(t_{0.025,\;n-1}\le T\le t_{0.975,\;n-1}\right)=0{,}95.
\]

Przypomnijmy, że $t_{p,\;n-1}$ oznacza \textbf{kwantyl rzędu $p$} rozkładu $t_{n-1}$, tzn.
\[
\mathbb{P}\!\left(T\le t_{0.025,\;n-1}\right)=0{,}025,
\qquad
\mathbb{P}\!\left(T\le t_{0.975,\;n-1}\right)=0{,}975.
\]
Stąd
\[
\mathbb{P}\!\left(T\ge t_{0.975,\;n-1}\right)=1-0{,}975=0{,}025.
\]
Zatem prawdopodobieństwo znalezienia się $T$ pomiędzy tymi kwantylami wynosi
\[
\mathbb{P}\!\left(t_{0.025,\;n-1}\le T\le t_{0.975,\;n-1}\right)
=1-0{,}025-0{,}025=0{,}95.
\]

Ponieważ rozkład $t_{n-1}$ jest symetryczny względem zera, kwantyle spełniają zależność
\[
t_{0.025,\;n-1}=-t_{0.975,\;n-1}.
\]
W konsekwencji równoważnie możemy zapisać
\[
\mathbb{P}\!\left(-t_{0.975,\;n-1}\le T\le t_{0.975,\;n-1}\right)=0{,}95.
\]

\smallskip
\textbf{Krok 2: podstawienie definicji $T$ i rozwiązanie nierówności względem $\mu$.}
Z kroku 2 mamy:
\[
\mathbb{P}\!\left(
-t_{0.975,\;n-1}
\le
\frac{\bar{X}-\mu}{S/\sqrt{n}}
\le
t_{0.975,\;n-1}
\right)=0{,}95.
\]
Mnożymy nierówność przez dodatnią liczbę $S/\sqrt{n}$ (nie zmienia znaku nierówności):
\[
\mathbb{P}\!\left(
-t_{0.975,\;n-1}\frac{S}{\sqrt{n}}
\le
\bar{X}-\mu
\le
t_{0.975,\;n-1}\frac{S}{\sqrt{n}}
\right)=0{,}95.
\]
Teraz przekształcamy tak, aby w środku był parametr $\mu$. 
\[
\mathbb{P}\!\left(
\bar{X}- t_{0.975,\;n-1}\frac{S}{\sqrt{n}}
\le
\mu
\le
\bar{X}+ t_{0.975,\;n-1}\frac{S}{\sqrt{n}}
\right)=0{,}95.
\]
Z definicji oznacza to, że
\[
\boxed{
CI_{0.95}(\mu)=
\left(
\bar{X}- t_{0.975,\;n-1}\frac{S}{\sqrt{n}},
\;
\bar{X}+ t_{0.975,\;n-1}\frac{S}{\sqrt{n}}
\right)
}
\]
jest $95\%$ przedziałem ufności dla $\mu$.

\end{itemize}


\item \textbf{Ocena interpretacji:}
\begin{quote}
``Z prawdopodobieństwem $95\%$ prawdziwa wartość średniej $\mu$ leży w wyznaczonym przedziale.''
\end{quote}

Ta interpretacja \textbf{nie jest poprawna} w sensie częstotliwościowym. Parametr $\mu$ jest stałą (nie jest zmienną losową), więc nie ma sensu mówić, że ``z prawdopodobieństwem $95\%$ $\mu$ leży w przedziale'' \emph{po} wyznaczeniu konkretnego przedziału na podstawie danych.

Poprawna interpretacja jest następująca:
gdybyśmy wielokrotnie losowali próby tej samej liczebności $n=14$ i za każdym razem konstruowali przedział ufności tą samą metodą, to \textbf{w $95\%$ takich prób} otrzymany przedział zawierałby prawdziwą wartość $\mu$, gdzie losowość dotyczy \textbf{przedziału} (bo zależy od losowej próby), a nie parametru $\mu$.

\end{enumerate}

\end{document}

